一单模无损腔的角频率为 \(\omega_c>0\)，一个二能级原子的跃迁角频率为 \(\omega_a>0\)，在旋转波近似下与腔模耦合，耦合频率 \(g>0\)。取 \(\Delta=\omega_a-\omega_c\)、\(\sigma_z=|e\rangle\langle e|-|g\rangle\langle g|\)，
\[
H=\hbar\omega_c a^\dagger a+\frac{\hbar\omega_a}{2}\sigma_z+
\hbar g(a\sigma_++a^\dagger\sigma_-).
\]
忽略腔损耗与自发辐射。必要时用 \(\bar n=|\alpha|^2\) 表示相干腔场平均光子数。
(1) 证明总激发数 \(N=a^\dagger a+\sigma_+\sigma_-\) 守恒。在 \(n\geq0\) 的 \(\{|e,n\rangle,|g,n+1\rangle\}\) 子空间精确对角化 \(H\)，求两条缀饰能级、本征态混合角，并说明 \(|g,0\rangle\) 的特殊地位。
(2) 原子初态 \(|e\rangle\)、腔场初态 \(|n\rangle\)。求任意失谐时的精确态矢、原子向基态跃迁概率及最大交换概率；分别核对光子数、总激发数和完整哈密顿量的能量守恒，讨论共振及大失谐极限。
(3) 令 \(\Delta=0\)，原子初态仍为 \(|e\rangle\)，但腔场改为相干态 \(|\alpha\rangle\)。从泊松光子数分布严格写出原子反转的级数；在 \(\bar n\gg1\) 的早期时间区间推导高斯塌缩包络与塌缩时间，并从相邻光子数的相位重合推导近似复苏时间，说明复苏不是严格周期。
(4) 在 \(|\Delta|\gg2g\sqrt{\bar n+1}\) 的色散区，从缀饰能级展开导出二阶有效哈密顿量（保留常数项以检验 \(|g,0\rangle\)）。若原子初态为 \((|e\rangle+|g\rangle)/\sqrt2\)，腔场为 \(|\alpha\rangle\)，求原子约化密度矩阵的非对角元及其模，解释腔场条件相位、可逆退相干和色散复相干时间。
